\documentclass[11pt]{article}

\usepackage[margin=1in]{geometry}
\usepackage{amsmath,amssymb,amsthm}
\usepackage{booktabs,array}
\usepackage{enumitem}
\usepackage{xcolor}
\usepackage[colorlinks=true,linkcolor=blue!50!black,citecolor=blue!50!black]{hyperref}

% Drafting helpers: delete calls as you fill in
\newcommand{\todo}[1]{\par\smallskip\noindent\textcolor{red!70!black}{\textbf{[FILL:}~#1\textbf{]}}\par\smallskip}
\newcommand{\decide}[1]{\par\smallskip\noindent\textcolor{orange!80!black}{\textbf{[DECIDE:}~#1\textbf{]}}\par\smallskip}

% Notation
\newcommand{\Gd}{G_{d}}          % those in power
\newcommand{\Gn}{G_{n}}          % general public
\newcommand{\Gr}{G_{r}}          % resistant subgroup (the movement)
\newcommand{\sd}{\sigma_{s}}     % dominant (silence) schema
\newcommand{\sr}{\sigma_{r}}     % resistant (disclosure) schema
\newcommand{\E}{\mathbb{E}}

\theoremstyle{definition}
\newtheorem{definition}{Definition}[section]
\theoremstyle{remark}
\newtheorem{remark}{Remark}[section]

\title{Formalizing Resistance}
\author{Aaron Avram}
\date{October 6, 2026}

\begin{document}
\maketitle

% ============================================================
\section{Set up}
% ============================================================
The basic setup for the system is as follows:
\begin{itemize}
    \item There is a group of agents $G$.
    \item A set of states $S$ where $s \in S$ has probability $p(s)$ of occuring.
    \item For each $s \in S$ there is a set $A(s)$ of possible actions in state $s$.
    \item A policy is function that maps each $s \in S$ to a probability distribution over $A(s)$.
    Fix a dominant policy $\pi_d$.
    \item Each agent $i \in G$ has a fixed reward $u_i(s, a)$ for action $a \in A(s)$ for $s \in S$.
    \item Each agent $i$ has an expected reward under a policy $\pi$ in state $s$:
    \[ U_i(s, \pi) = \mathbb E_{a \sim \pi(s)}\Big[u_i(s, a)\Big] \]
    \item For each $i \in G$ there is a set of states $S_r^i \subseteq S$ where $U_i(s, \pi_d) \ll 0$.
\end{itemize}
We will assume the dominant norm stays fixed throughout all of the calculations for simplicity.

\section{Defining Resistance Theoretically}
We give a theoretical definition of a resistant subgroup $G_r$ of $G$.
\[ w_i(s) = \max(0, -U_i(s, \pi_d)) \]
hence $w_i \in \mathbb R^{|S|}_{\geq 0}$.
The probability function on $S$ defines a norm on $\mathbb R^{|S|}$ given by:
\[ \| \cdot \|_p: x \to \sqrt{\sum_{s \in S}p(s)x(s)^2} \]
For a subgroup $H \subseteq G$ define
\[ D(H) = \sum_{i \in H} \| w_i \|_p^2 \]
which is the total disagreement with $\pi_d$ within $H$.
If $\overline{w}_H$ is the mean $w_i$ in $H$, we define a related quantity:
\[ F(H) = \sum_{i \in H}\| w_i -\overline{w}_H \|_p^2 \]
and by expanding the square this is equivalent to
\[ F(H) = \frac{1}{2|H|}\sum_{i, j \in H}\| w_i - w_j \|^2_p \]
this is called the friction of $H$ and represents how much cohesion the group $H$
has in its disagreement with $\pi_d$.

Given a tolerance $\delta > 0$ we define the maximally resistant subgroup of $G$
as
\[ G_r = \operatorname{argmax}_{H \in \mathcal P(G)}D(H) \]
subject to the constraint
\[ F(H)/|H| \leq \delta \]
where the constraint is necessary as $D(H)$ is monotone in $H$ with respect to $\subseteq$.

If $\delta = 0$ then this requires that the maximally resistant subgroup is frictionless, i.e.
its members perfectly align on their disagreement with $\pi_d$.

Now we define a similarity matrix for the set of agents.
\[ K_{ij} = \exp(- \| w_i - w_j \|^2 /2\sigma^2) \]
and note that Jensen's inequality gives the relation:
\[ \mathbb E_H[K_{ij}] \leq \exp \mathbb E_H[\| w_i - w_j \|^2_p / 2\sigma^2] =  \exp (\frac{F(H)}{|H|\sigma^2}) \]
which gives the bound
\[ |H|\sigma^2 \log \mathbb E_H[K_{ij}] \leq F(H) \]
To prove a flipped version of this inequality we assume that $|u_i(s, a)|$ is bounded by a uniform constant.
This implies that there exists $R$ such that for all $i, j \in H$
\[ \| w_i - w_j \|_p^2 \leq R \]
let $M = \frac{R^2}{2\sigma^2}$. Then the function
\[ 1 - \exp(-x) \]
is concave so for $x \in [0, M]$
\[ 1 - \exp(-x) \geq x\Big(1 - \exp(-M)\Big)/M \]
rearranging we see that
\[ x \leq M\frac{1 - \exp(-x)}{1 - \exp(-M)} \]
now set $x = \|w_i - w_j \|_p^2/2\sigma^2$
then
\[ \|w_i - w_j \|_p^2/2\sigma^2 \leq \frac{R^2}{2\sigma^2}\frac{1 - K_{ij}}{1 - \exp(-M)} \]
or
\[ \|w_i - w_j \|_p^2 \leq R^2\frac{1 - K_{ij}}{1 - \exp(-M)} \]
hence
\[ F(H) \leq \frac{R^2}{2}\frac{\frac{1}{|H|} - \mathbb E[K_{ij}]}{1 - \exp(-\frac{R^2}{2\sigma^2})} \]
in other words this implies that the expected pairwise distance is small whenever $F(H)$ is.

Since $F(H)$ would be hard to compute in practice $K_{ij}$ can be used as an approximation.

\section{Estimating Rewards}
Each agent keeps a running average of their personal benefit under $\pi_d$
(this defined similarly to the learning norms paper). Now I propose that each agent
shares their estimates with all other agents. This forms a matrix:
\[ V_i[j, s] \]
which tracks $i$'s estimate of $j$'s personal benefit under the dominant norm in state $s$.
$V_i[j, s]$ should approach $U_j(s, \pi_d)$.

Currently I assume that the message passed from agent $j$ to $i$ is perfectly accurate.
However, this could be extended in two ways: one, simulating imperfect communication
by adding Gaussian noise to each message and two, simulating deception by allowing $j$
to learn the message to pass to $i$ that is most favourable to it.

Next we define
\[ \hat{w}_j^i(s) = \max(0, -V_i[j, s]) \]
under the above assumptions this is independent of $i$ so we just write $\hat{w}_j(s)$.
Additionally define
\[ \hat{d}^2_{jk} = \| \hat{w}_j(s) - \hat{w}_k(s) \|^2_p \]
and
\[ \hat{K}_{jk} = \exp(-\hat{d}^2_{jk}/\sigma^2) \]
then
\[ \epsilon_j = \| w_j - \hat{w}_j \|_p \]
with rewards bounded in absolute value by $B$ Hoeffding's inequality
gives the bound
\[ | w_j(s) - \hat{w}_j(s) | < \sqrt{2\log(2/\delta')/n_j(s)} \]
with probability $1 - \delta'$ where $n_j(s)$ is the number of times $j$ has been in state $s$
so far in the system. This then gives the bound
\[ \epsilon_j < \sqrt{2|S|\log(2|S||N|/\delta')/n_j} \]
where $n_j$ is the number of times agent $j$ has acted so far.

This shows that $\epsilon_j$ can be bouned by quantities observable by the agents,
but for simplicity I will assume below that each agent $j$ has access to $\epsilon_j$.

\section{Formation}
We now define conditions for resistant subgroup formation.

An agent $i$ is called a candidate if
\[ \| w_i \|^2_p \geq \tau \]
for some fixed constant $\tau$.

Then $i$ admits $j$ to its group if
\[ \hat{d}_{ij} - \epsilon_i - \epsilon_j < r_{in} \]
and $i$ leaves if some member $j$ has
\[ \hat{d}_{ij} - \epsilon_i - \epsilon_j > r_{out} \]
where $r_{out} > r_{in}$ for hysteresis.

These kinds of decisions would occur at update epochs like in the Learning Norms paper.

\end{document}